\begin{abstract}
A meeting interpreter must say what each person said, attribute it to the right person, and get the translation to
the listener in time, also when several people talk at once. In \sys{}, a speech LLM reads acoustic tokens interleaved
with tokens of a frozen streaming diarizer, writes speaker-first, time-stamped translations for every speaker of an
overlapping window, and hands them to a speech generator prompted with each speaker's voice. Offline, \sys{} beats the
strongest non-oracle diarize-then-translate cascade by \num{v7:gainoff:zh2en} cpBLEU in Chinese$\rightarrow$English
(confirmatory, Holm-corrected), and on overlapped speech it leads a speaker-blind model relabelled by the same diarizer,
mainly by attributing both voices. Online, much of the delay a listener experiences arises after the text is written,
while synthesised speech queues behind earlier speech of the same speaker. \emph{Playback alignment} scores each
block's candidate renderings together with this backlog in a linear-chain CRF whose potentials charge the delay at
which each word is heard, imposes a latency budget by posterior regularisation, and fine-tunes the speech generator once on the posterior-weighted candidates. On the same
text, on whole AMI meetings (Zh$\rightarrow$En / En$\rightarrow$Zh) and real Mandarin AliMeeting meetings, it raises
speech ASR-BLEU by \num{v10:sig:test:ami:speech:e2-b:diff} / \num{v10:sig:test:en:speech:e2-b:diff} /
\num{v10:sig:test:ali:speech:e2-b:diff}, lowers conventional block-start speech LAAL by
\num{v10:pr:test:ami:e2-b:bslaal:diffabs} / \num{v10:pr:test:en:e2-b:bslaal:diffabs} /
\num{v10:pr:test:ali:e2-b:bslaal:diffabs}\,s and ends each utterance's speech \num{v10:pr:test:ami:e2-b:endoff:diffabs} /
\num{v10:pr:test:en:e2-b:endoff:diffabs} / \num{v10:pr:test:ali:e2-b:endoff:diffabs}\,s earlier (all significant), by speaking faster with less queueing and at the price of about twice as many
pauses.\ifFactorial{} A factorial ablation separates the learned component from merging, rate rules and candidate
selection.\fi\ifAliEval{} \num{v11:ae:abstract}\fi{} Treating the write time as latent structure instead, trained on the decoder's own likelihood, brings no
significant gain: its posterior drifts towards later writes, and we analyse why neither a training-time budget, an
early-write lattice nor a decoding bias turns it into a better quality--latency trade-off.
\end{abstract}
